\documentclass{article}
\usepackage{../assignment_preamble}

\title{Weekly Problem 2}
\author{Ravi Kini}
\date{April 13, 2026}

\begin{document}

\maketitle

\problem
\subproblem{(a)}
\begin{equation*}
\begin{nd}[rules=strongkleene,justformat=just]
\hypo {1} {\lnot p}
\open
\hypo {2} {p \land q}
\have {3} {p} \ae{2}
\have {4} {\lnot\lnot p} \dn{3}
\have {5} {\lnot\lnot p \lor q} \oi{4}
\have {6} {\lnot p \Rightarrow q} \mci{5}
\close
\have {7} {(p \land q) \Rightarrow (\lnot p \Rightarrow q)} \ii{2-6}
\end{nd}
\end{equation*}
We then see that $\lnot p \vdash (p \land q) \Rightarrow (\lnot p \Rightarrow q)$ is a valid argument in Strong Kleene logic.

\subproblem{(b)}
\begin{equation*}
\begin{nd}[rules=strongkleene,justformat=just]
\hypo {1} {p \Rightarrow q}
\hypo {2} {\lnot p \Rightarrow q}
\hypo {3} {q \lor \lnot q}
\open
\hypo {4} {q}
\have {5} {q} \r{4}
\close
\open
\hypo {5} {\lnot q}
\have {6} {\lnot p \lor q} \mce{1}
\open
\hypo {7} {\lnot p}
\have {8} {q} \ie{2,7}
\close
\open
\hypo {9} {q}
\have {10} {q} \r{10}
\close
\have {11} {q} \oe{6,7-8,9-10} 
\close
\have {11} {q} \oe{3,4-5,6-11} 
\end{nd}
\end{equation*}
We then see that $p \Rightarrow q, \lnot p \Rightarrow q, q \lor \lnot q \vdash q$ is a valid argument in Strong Kleene logic.

\end{document}